\documentclass[10pt,a4paper]{article}
\usepackage[brazil]{babel}
\usepackage[utf8]{inputenc}
\renewcommand{\baselinestretch}{1.0}
\usepackage{mhchem}
\usepackage{graphicx}
\usepackage{amssymb}
\usepackage{amsmath}
\usepackage{enumitem}
\usepackage{geometry}
\geometry{a4paper, margin=2cm}
\setlist{nosep, topsep=2pt, partopsep=0pt, parsep=0pt, itemsep=3pt}

\newcommand{\oC}{$^\circ$C}
\newcommand{\e}[1]{$\times 10^{#1}$}
\sffamily

\newcommand{\questaoheader}[1]{%
  \begin{center}
  {\small\sffamily QG101 --- Bloco 13: Termoquímica Básica
  \hfill Questão #1}
  \end{center}
  \vspace{-0.3cm}
  \noindent\rule{\textwidth}{0.6pt}
  \vspace{0.1cm}
}

\newcommand{\resolucaosep}{%
  \vspace{0.4cm}
  \noindent\rule{\textwidth}{0.4pt}
  \vspace{0.1cm}

  \noindent{\bf Resolução:}
  \vspace{0.2cm}
}

\begin{document}
\pagestyle{empty}

% ================================================================
% Questão 1
% ================================================================
\questaoheader{1}

\noindent\textbf{Questão 1.}
Classifique cada processo abaixo como \textbf{exotérmico} ou
\textbf{endotérmico}, justificando em uma frase:
\begin{enumerate}
  \item Combustão de madeira.
  \item Derretimento de gelo.
  \item Fotossíntese.
  \item Respiração celular.
\end{enumerate}

\resolucaosep

\begin{enumerate}
  \item \textbf{Exotérmico}: a combustão oxida a celulose e outros
        compostos da madeira, liberando calor e luz para o ambiente
        ($\Delta H < 0$).

  \item \textbf{Endotérmico}: o gelo absorve calor do ambiente para romper
        as ligações de hidrogênio da estrutura cristalina e se tornar líquido
        ($\Delta H > 0$).

  \item \textbf{Endotérmico}: a fotossíntese usa energia luminosa (solar)
        para converter \ce{CO2} e \ce{H2O} em glicose e \ce{O2} ---
        armazena energia no sistema ($\Delta H > 0$).

  \item \textbf{Exotérmico}: a respiração celular oxida a glicose,
        liberando energia química na forma de ATP e calor para o organismo
        ($\Delta H < 0$).
\end{enumerate}

\clearpage

% ================================================================
% Questão 2
% ================================================================
\questaoheader{2}

\noindent\textbf{Questão 2.}
Uma reação química libera $50$~kJ para a vizinhança.
\begin{enumerate}
  \item Qual é o sinal de $\Delta H$ dessa reação?
  \item A reação é exotérmica ou endotérmica? Justifique.
  \item Esboce, em palavras, como seria o diagrama de energia
        (reagentes x produtos) dessa reação.
\end{enumerate}

\resolucaosep

\begin{enumerate}
  \item $\Delta H = -50$~kJ (sinal \textbf{negativo}): por convenção,
        $\Delta H = H_{\text{produtos}} - H_{\text{reagentes}}$; como o
        sistema perde energia para a vizinhança, os produtos ficam num
        nível de entalpia mais baixo que os reagentes.

  \item A reação é \textbf{exotérmica}: energia é liberada pelo sistema
        para a vizinhança ($\Delta H < 0$), e a vizinhança aquece.

  \item O diagrama de energia (entalpia vs.\ coordenada de reação) mostra
        os \textbf{reagentes num patamar mais alto}, uma barreira de
        ativação acima deles (energia de ativação), e os \textbf{produtos
        num patamar mais baixo}, com a diferença de altura correspondendo
        a 50~kJ --- a energia liberada na forma de calor.
\end{enumerate}

\clearpage

% ================================================================
% Questão 3
% ================================================================
\questaoheader{3}

\noindent\textbf{Questão 3.}
Calcule a quantidade de calor necessária para aquecer $100$~g de água de
$25$\oC{} a $75$\oC{} (calor específico da água
$= 4{,}18$~J~g$^{-1}$\oC$^{-1}$). Expresse o resultado em joules e em
quilojoules.

\resolucaosep

Usando $q = m \cdot c \cdot \Delta T$:

\begin{align*}
\Delta T &= 75 - 25 = 50\text{\oC}\\
q &= 100\,\text{g} \times 4{,}18\,\text{J\,g}^{-1}\text{\oC}^{-1}
     \times 50\text{\oC}\\
q &= \mathbf{20\,900\,\text{J}} = \mathbf{20{,}9\,\text{kJ}}
\end{align*}

O calor necessário é \textbf{positivo} (o sistema absorve energia), o que
é coerente com um processo de aquecimento.

\clearpage

% ================================================================
% Questão 4
% ================================================================
\questaoheader{4}

\noindent\textbf{Questão 4.}
Dada a equação termoquímica:
$$\ce{C(s) + O2(g) -> CO2(g)} \qquad \Delta H = -394~\text{kJ/mol}$$
\begin{enumerate}
  \item O que significa o sinal negativo de $\Delta H$?
  \item Quantos kJ são liberados na combustão completa de $2{,}00$~mol de
        \ce{C(s)}?
  \item Quantos kJ são liberados na combustão completa de $6{,}0$~g de
        \ce{C(s)} ($M = 12{,}0$~g/mol)?
\end{enumerate}

\resolucaosep

\begin{enumerate}
  \item O sinal negativo indica que a reação é \textbf{exotérmica}: o
        sistema (carbono + oxigênio) libera 394~kJ de calor para a
        vizinhança a cada mol de \ce{C} consumido. Os produtos possuem
        entalpia menor que os reagentes.

  \item
  \[
    q = 2{,}00\,\text{mol} \times 394\,\text{kJ/mol}
      = \mathbf{788\,\text{kJ}}
  \]
  São liberados 788~kJ.

  \item Número de mols de \ce{C}:
  \[
    n = \frac{6{,}0\,\text{g}}{12{,}0\,\text{g/mol}} = 0{,}50\,\text{mol}
  \]
  \[
    q = 0{,}50\,\text{mol} \times 394\,\text{kJ/mol}
      = \mathbf{197\,\text{kJ}}
  \]
  São liberados 197~kJ.
\end{enumerate}

\clearpage

% ================================================================
% Questão 5
% ================================================================
\questaoheader{5}

\noindent\textbf{Questão 5.}
Por que a transpiração resfria o corpo? Explique em termos de processo
endotérmico/exotérmico e do tipo de mudança de estado envolvida.

\resolucaosep

A transpiração envolve a \textbf{vaporização} (evaporação) do suor ---
passagem do estado \textbf{líquido para o gasoso}.

A vaporização é um processo \textbf{endotérmico}: as moléculas de água
precisam absorver energia para vencer as forças intermoleculares
(ligações de hidrogênio) que as mantêm no líquido e escapar para a fase
de vapor. Essa energia é retirada da \emph{superfície da pele} e dos
tecidos próximos, que se resfriam.

Em termos quantitativos, a entalpia de vaporização da água é elevada
($\approx 2260$~J/g a 37\oC), o que torna a transpiração um mecanismo
de termorregulação muito eficiente: basta evaporar pequenas quantidades
de suor para dissipar grandes quantidades de calor corporal.

\clearpage

% ================================================================
% Questão 6
% ================================================================
\questaoheader{6}

\noindent\textbf{Questão 6.}
São conhecidos os seguintes dados termoquímicos:
\begin{align*}
\ce{H2(g) + 1/2 O2(g) &-> H2O(l)} &\Delta H_1 &= -286~\text{kJ/mol}\\
\ce{H2(g) + 1/2 O2(g) &-> H2O(g)} &\Delta H_2 &= -242~\text{kJ/mol}
\end{align*}
\begin{enumerate}
  \item Usando a Lei de Hess, calcule o $\Delta H$ da vaporização da água,
        $\ce{H2O(l) -> H2O(g)}$.
  \item O valor obtido deveria ser positivo ou negativo? Esse resultado é
        coerente com o tipo de processo (fornecimento ou liberação de
        calor)?
\end{enumerate}

\resolucaosep

\begin{enumerate}
  \item Pela Lei de Hess, combinamos as equações de modo que os termos
        indesejados se cancelem. Invertemos a equação 1 (mudando o sinal
        de $\Delta H_1$) e somamos com a equação 2:

        \begin{align*}
          \text{(1 invertida):}\quad
          \ce{H2O(l) &-> H2(g) + 1/2 O2(g)}
          &\Delta H &= +286\,\text{kJ/mol}\\
          \text{(2):}\quad
          \ce{H2(g) + 1/2 O2(g) &-> H2O(g)}
          &\Delta H &= -242\,\text{kJ/mol}
        \end{align*}

        Somando:
        \[
          \ce{H2O(l) -> H2O(g)}
          \qquad
          \Delta H_{\text{vap}} = 286 + (-242) = \mathbf{+44\,\text{kJ/mol}}
        \]

  \item O valor é \textbf{positivo}, o que é coerente: a vaporização é um
        processo \textbf{endotérmico} --- o sistema (água) precisa absorver
        energia para romper as ligações de hidrogênio e passar para o
        estado gasoso.
\end{enumerate}

\clearpage

% ================================================================
% Questão 7
% ================================================================
\questaoheader{7}

\noindent\textbf{Questão 7.}
Uma amostra de $100$~g de água e uma amostra de $100$~g de ferro
($c = 0{,}45$~J~g$^{-1}$\oC$^{-1}$), inicialmente à mesma temperatura,
recebem a mesma quantidade de calor ($5{,}0$~kJ cada).
\begin{enumerate}
  \item Calcule a variação de temperatura ($\Delta T$) de cada amostra.
  \item Qual delas esquenta mais? Relacione o resultado ao conceito de
        calor específico e a um exemplo do cotidiano (ex.: areia da praia
        x água do mar).
\end{enumerate}

\resolucaosep

Usando $q = m \cdot c \cdot \Delta T \;\Rightarrow\; \Delta T = q/(m \cdot c)$:

\begin{enumerate}
  \item
    \begin{itemize}
      \item \textbf{Água} ($c = 4{,}18$~J~g$^{-1}$\oC$^{-1}$):
        \[
          \Delta T_{\text{água}}
          = \frac{5000\,\text{J}}{100\,\text{g}
            \times 4{,}18\,\text{J\,g}^{-1}\text{\oC}^{-1}}
          = \frac{5000}{418}
          \approx \mathbf{12{,}0\text{\oC}}
        \]

      \item \textbf{Ferro} ($c = 0{,}45$~J~g$^{-1}$\oC$^{-1}$):
        \[
          \Delta T_{\text{ferro}}
          = \frac{5000\,\text{J}}{100\,\text{g}
            \times 0{,}45\,\text{J\,g}^{-1}\text{\oC}^{-1}}
          = \frac{5000}{45}
          \approx \mathbf{111\text{\oC}}
        \]
    \end{itemize}

  \item O \textbf{ferro} esquenta muito mais (111\oC vs.\ 12\oC) porque
        seu calor específico é quase 10 vezes menor que o da água. Um
        calor específico baixo significa que pouca energia é necessária
        para elevar a temperatura de 1\,g em 1\oC.

        Exemplo do cotidiano: a \textbf{areia da praia} (baixo $c$)
        esquenta rapidamente ao sol e fica quase insuportável ao toque,
        enquanto a \textbf{água do mar} (alto $c$) permanece bem mais
        fria --- ela absorve muito mais calor para a mesma variação de
        temperatura.
\end{enumerate}

\clearpage

% ================================================================
% Questão 8
% ================================================================
\questaoheader{8}

\noindent\textbf{Questão 8.}
A síntese da amônia é descrita por:
$$\ce{N2(g) + 3H2(g) -> 2NH3(g)} \qquad \Delta H = -92~\text{kJ/mol}$$
\begin{enumerate}
  \item Qual o valor de $\Delta H$ para a formação de \textbf{1 mol} de
        \ce{NH3(g)} a partir de \ce{N2} e \ce{H2}?
  \item Escreva a equação termoquímica da reação \textbf{inversa}
        (decomposição de $2$~mol de \ce{NH3}) e indique seu $\Delta H$.
\end{enumerate}

\resolucaosep

\begin{enumerate}
  \item A equação fornecida produz 2~mol de \ce{NH3} com $\Delta H =
        -92$~kJ. Para 1~mol de \ce{NH3}, divide-se tudo por 2 (a
        estequiometria e o $\Delta H$):

        \[
          \tfrac{1}{2}\ce{N2(g) + \tfrac{3}{2}H2(g) -> NH3(g)}
          \qquad
          \Delta H = \frac{-92}{2} = \mathbf{-46\,\text{kJ/mol}}
        \]

  \item Invertendo a reação original (sentido inverso), o sinal de
        $\Delta H$ também se inverte:

        \[
          \ce{2NH3(g) -> N2(g) + 3H2(g)}
          \qquad
          \Delta H = \mathbf{+92\,\text{kJ/mol}}
        \]

        A decomposição da amônia é \textbf{endotérmica}: requer
        fornecimento de energia para romper as ligações N--H e reforjar
        a tripla ligação N$\equiv$N e as ligações H--H.
\end{enumerate}

\clearpage

% ================================================================
% Questão 9
% ================================================================
\questaoheader{9}

\noindent\textbf{Questão 9.}
Um gás, ao absorver $300$~J de calor da vizinhança, se expande e realiza
$120$~J de trabalho sobre a vizinhança.
\begin{enumerate}
  \item Identifique o sinal de $q$ e de $w$ nesse processo, usando a
        convenção $\Delta U = q + w$ (em que $w$ é o trabalho realizado
        \emph{sobre} o sistema).
  \item Calcule $\Delta U$ do gás.
\end{enumerate}

\resolucaosep

\begin{enumerate}
  \item Na convenção $\Delta U = q + w$ (convenção IUPAC/química):
        \begin{itemize}
          \item $q$: calor \emph{absorvido pelo sistema} → \textbf{positivo}.
                O gás absorbe 300~J $\Rightarrow q = +300$~J.

          \item $w$: trabalho realizado \emph{sobre o sistema}. O sistema
                (gás) se \emph{expande} e realiza trabalho \emph{sobre a
                vizinhança} (trabalho feito pelo sistema), portanto o
                trabalho feito \emph{sobre o sistema} é negativo:
                $w = -120$~J.
        \end{itemize}

  \item
        \[
          \Delta U = q + w = 300 + (-120) = \mathbf{+180\,\text{J}}
        \]
        A energia interna do gás \textbf{aumentou} em 180~J: o sistema
        ganhou mais energia via calor do que cedeu via trabalho de
        expansão.
\end{enumerate}

\clearpage

% ================================================================
% Questão 10
% ================================================================
\questaoheader{10}

\noindent\textbf{Questão 10.}
\textbf{Calorimetria.} Um pedaço de $50{,}0$~g de um metal, inicialmente a
$100{,}0$\oC, é colocado em $100{,}0$~g de água a $20{,}0$\oC. O sistema
atinge o equilíbrio térmico a $24{,}0$\oC. Use
$c_{\text{água}} = 4{,}18$~J~g$^{-1}$\oC$^{-1}$ e admita que todo o calor
perdido pelo metal é ganho pela água (sistema isolado).
\begin{enumerate}
  \item Calcule o calor ganho pela água.
  \item Calcule o calor específico do metal, sabendo que ele perde a
        mesma quantidade de calor que a água ganha.
\end{enumerate}

\resolucaosep

\begin{enumerate}
  \item \textbf{Calor ganho pela água} ($\Delta T_{\text{água}} =
        24{,}0 - 20{,}0 = 4{,}0$\oC):
        \[
          q_{\text{água}}
          = m \cdot c \cdot \Delta T
          = 100{,}0\,\text{g}
            \times 4{,}18\,\text{J\,g}^{-1}\text{\oC}^{-1}
            \times 4{,}0\text{\oC}
          = \mathbf{1672\,\text{J}}
        \]

  \item O metal perde $q_{\text{metal}} = -1672$~J (sinal negativo:
        o metal cede calor).

        $\Delta T_{\text{metal}} = 24{,}0 - 100{,}0 = -76{,}0$\oC

        Isolando $c_{\text{metal}}$:
        \[
          q_{\text{metal}}
            = m_{\text{metal}} \cdot c_{\text{metal}} \cdot \Delta T_{\text{metal}}
        \]
        \[
          -1672
            = 50{,}0 \times c_{\text{metal}} \times (-76{,}0)
        \]
        \[
          c_{\text{metal}}
            = \frac{1672}{50{,}0 \times 76{,}0}
            = \frac{1672}{3800}
            \approx \mathbf{0{,}44\,\text{J\,g}^{-1}\text{\oC}^{-1}}
        \]

        Esse valor é muito próximo do calor específico do \textbf{ferro}
        ($c = 0{,}45$~J~g$^{-1}$\oC$^{-1}$), sugerindo que o metal em
        questão pode ser o ferro.
\end{enumerate}

\end{document}
